\documentclass[10pt,twocolumn]{article}
\usepackage[T1]{fontenc}
\usepackage[utf8]{inputenc}
\usepackage{lmodern}
\usepackage[a4paper,top=22mm,bottom=24mm,left=18mm,right=18mm,columnsep=6mm]{geometry}
\usepackage{graphicx}
\usepackage{amsmath}
\usepackage{booktabs}
\usepackage{microtype}
\usepackage[hidelinks]{hyperref}
\usepackage{xcolor}
\usepackage{listings}
\usepackage{titlesec}
\usepackage{fancyhdr}
\usepackage{lastpage}
\usepackage{natbib}
\bibpunct{(}{)}{;}{a}{}{,}

\definecolor{jpteal}{HTML}{0E8F9E}
\titleformat{\section}{\normalfont\large\bfseries}{\thesection}{0.6em}{}
\titleformat{\subsection}{\normalfont\normalsize\bfseries}{\thesubsection}{0.6em}{}
\titlespacing*{\section}{0pt}{1.4ex plus .3ex}{0.7ex}
\titlespacing*{\subsection}{0pt}{1.1ex plus .2ex}{0.5ex}
\setlength{\parskip}{0pt}
\setlength{\parindent}{1.2em}
\raggedbottom
\clubpenalty=10000
\widowpenalty=10000
\pagestyle{fancy}\fancyhf{}
\renewcommand{\headrulewidth}{0.4pt}
\lhead{\small\textit{JuPIPER}}
\rhead{\small\thepage\ of \pageref{LastPage}}

\newcommand{\code}[1]{\texttt{\small #1}}

% Every campaign and real-complex figure below is a macro computed from the
% committed result files by make_numbers.py. Regenerate after new data; do not
% type a number from either file into the prose.
\graphicspath{{build/}}
% GENERATED by reproducibility/make_numbers.py -- do not edit.
% Sources:
%   bench/results/wp8-batch-graniterapids.json
%   bench/results/wp8-batch-tremont-20260820t105514z.json
%   bench/results/realdock-1ggr-f32-tremont-r5-nvidia-geforce-rtx-4090-gpu.json
%   bench/results/realdock-1ggr-f64-tremont-r5-nvidia-geforce-rtx-4090-gpu.json
%   results/wp8/comparison_summary.csv

% ---- the benchmark campaign, first run ----
\newcommand{\campCases}{52}
\newcommand{\campOk}{51}
\newcommand{\campRankOne}{27}
\newcommand{\campRankOnePct}{52.9}
\newcommand{\campHitTopK}{44}
\newcommand{\campHitTopKPct}{86.3}
\newcommand{\campClusterOne}{15}
\newcommand{\campClusterOnePct}{29.4}
\newcommand{\campIrmsdFour}{27}
\newcommand{\campSeconds}{1909.0}
\newcommand{\campMinutes}{31.8}
\newcommand{\campPerCaseLo}{34.8}
\newcommand{\campPerCaseHi}{45.1}
\newcommand{\campFloorLo}{0.42}
\newcommand{\campFloorHi}{1.20}
\newcommand{\campFloorMed}{0.73}
\newcommand{\campAngleFloor}{2.37}
\newcommand{\campNrot}{69{,}960}
\newcommand{\campWone}{7}
\newcommand{\campHitAll}{50}
\newcommand{\campMaxHitRank}{24{,}544}
\newcommand{\campEqualRanks}{51}
\newcommand{\campEqualEnergy}{35}
\newcommand{\dockBestRank}{293}
\newcommand{\campTableBody}{1ACB & 1 & 206 & 1 & 40.9 \\
1AVX & 19 & 42 & 7 & 40.8 \\
1AY7 & 490 & 2 & 26 & 34.9 \\
1BVN & 1 & 351 & 1 & 36.9 \\
1CGI & 1 & 567 & 1 & 38.6 \\
1CLV & 1 & 465 & 1 & 35.2 \\
1D6R & -- & 0 & -- & 36.0 \\
1DFJ & 1 & 16 & 17 & 35.5 \\
1E6E & 3 & 139 & 2 & 34.8 \\
1EAW & 95 & 9 & 11 & 37.3 \\
1EWY & 204 & 23 & 8 & 35.8 \\
1EZU & 3 & 67 & 3 & 39.1 \\
1F34 & 1 & 130 & 2 & 36.0 \\
1F6M & 9 & 64 & 5 & 39.6 \\
1FLE & 1 & 255 & 1 & 36.9 \\
1FQ1 & 4938 & 0 & -- & 36.5 \\
1GL1 & 1 & 175 & 2 & 35.6 \\
1GXD & 25 & 31 & 9 & 36.3 \\
1HIA & 1 & 186 & 1 & 36.4 \\
1IJK & 121 & 16 & 12 & 35.8 \\
1JIW & 22 & 8 & 21 & 37.8 \\
1JTG & 1 & 76 & 2 & 36.3 \\
1KKL & 1615 & 0 & -- & 36.5 \\
1M10 & 451 & 1 & -- & 38.6 \\
1MAH & 1 & 78 & 6 & 41.3 \\
1N8O & 1 & 150 & 2 & 38.6 \\
1NW9 & 1 & 215 & 1 & 41.0 \\
1OC0 & 2 & 187 & 11 & 40.6 \\
1OPH & 19 & 45 & 6 & 40.3 \\
1OYV\_A & 1 & 134 & 2 & 36.5 \\
1OYV\_B & 5130 & 0 & -- & 45.1 \\
1PPE & 1 & 172 & 3 & 38.5 \\
1PXV & 1 & 31 & 11 & 37.0 \\
1R0R & 1 & 147 & 4 & 36.2 \\
1TMQ & 1 & 209 & 2 & 35.5 \\
1UDI & 1 & 245 & 1 & 38.6 \\
1YVB & 1 & 151 & 1 & 36.4 \\
1ZLI & 1 & 174 & 1 & 36.4 \\
2ABZ & 19 & 105 & 4 & 39.1 \\
2B42 & 1 & 39 & 7 & 35.7 \\
2J0T & 1 & 152 & 1 & 38.6 \\
2MTA & 4 & 254 & 1 & 35.6 \\
2O3B & 2948 & 0 & -- & 35.8 \\
2O8V & 1 & 196 & 2 & 35.6 \\
2OUL & 1 & 163 & 1 & 36.2 \\
2PCC & 5720 & 0 & -- & 39.0 \\
2SIC & 1 & 279 & 1 & 38.6 \\
2SNI & 1 & 211 & 1 & 35.9 \\
2UUY & 927 & 1 & -- & 37.5 \\
3SGQ & 376 & 8 & 10 & 35.5 \\
4CPA & \multicolumn{4}{l}{Atom typing failed} \\
7CEI & 24544 & 0 & -- & 35.8 \\}

% ---- the second campaign, a different machine ----
\newcommand{\campTwoSeconds}{2563.5}
\newcommand{\campTwoMinutes}{42.7}
\newcommand{\campTwoRankOne}{27}
\newcommand{\campTwoHitTopK}{44}
\newcommand{\campTwoClusterOne}{15}

% ---- the real complex, both arms in one process ----
\newcommand{\dockGpuFThirtyTwo}{20.7}
\newcommand{\dockGpuFSixtyFour}{60.5}
\newcommand{\dockCpuFThirtyTwo}{1083.7}
\newcommand{\dockSpeedupFThirtyTwo}{52.3}
\newcommand{\dockSpeedupFSixtyFour}{18.7}
\newcommand{\dockAtoms}{1{,}774}
\newcommand{\dockNrot}{69{,}960}
\newcommand{\dockRank}{1}
\newcommand{\dockRmsd}{8.911}
\newcommand{\dockBest}{0.85974}
\newcommand{\dockFloor}{0.85974}
\newcommand{\dockEnergy}{-659.041}
\newcommand{\parityCompared}{2{,}000}
\newcommand{\parityCellAgree}{1.0}
\newcommand{\parityRelEnergy}{5.9e-06}

% ---- the external comparison, five cases ----
\newcommand{\wpCases}{5}
\newcommand{\wpJoint}{0}
\newcommand{\wpSameBand}{3}
\newcommand{\wpNearN}{4}
\newcommand{\wpNearLo}{1.85}
\newcommand{\wpNearHi}{1.98}
\newcommand{\wpNearSpread}{0.13}
\newcommand{\wpRankLo}{20}
\newcommand{\wpRankHi}{437}
\newcommand{\wpCentreLo}{3.77}
\newcommand{\wpCentreHi}{6.20}
\newcommand{\wpCentreAllLo}{3.77}
\newcommand{\wpCentreAllHi}{6.68}
\newcommand{\wpOddCase}{1AY7}
\newcommand{\wpOddNear}{5.76}
\newcommand{\wpTableBody}{1ACB & 1 & 1 & 4.58 & 1 & 1.93 & 319 \\ 1AVX & 1 & 7 & 5.60 & 7 & 1.85 & 20 \\ 1AY7 & -- & 26 & 6.68 & 7 & 5.76 & 338 \\ 1BVN & 1 & 1 & 3.77 & 1 & 1.98 & 106 \\ 1CGI & 1 & 1 & 6.20 & 1 & 1.98 & 437 \\}

% ---- carried by hand: no committed JSON holds these ----
%   tests on CPU                       5,205                          <- console output of Pkg.test()
%   tests with the GPU extension       5,282                          <- console output of a device deploy
%   correlation vs direct O(N^6)       5.1e-16                        <- test_correctness.jl
%   CPU/GPU agreement, Float32         3e-7                           <- device suite, paper1_hpc 6.5.1
%   CPU/GPU agreement, Float64         7e-16                          <- device suite, paper1_hpc 6.5.1
%   pipeline vs its own floor, CPU     1.11-1.52x                     <- paper1_hpc 6.3
%   pipeline vs its own floor, GPU     1.17-2.50x                     <- paper1_hpc 6.5.2
%   Fourier-domain accumulation        1.56-1.97x                     <- paper1_hpc 6.4 row A1
%   real-input transforms              1.44-3.29x                     <- paper1_hpc 6.4 row A2
%   translation scan rewrite           16.5-60.3x                     <- paper1_hpc 6.5.5
%   closed-form pose RMSD              1225x at M=2000                <- paper1_hpc 2.6
%   shape-weight crossover             3.264 / 3.151                  <- paper1_hpc 6.8
%   fibre pruning                      0 of 200, 1.34x slower         <- paper1_hpc 7.1
%   restraints                         78.06% infeasible, 1.26-1.86x  <- paper1_hpc 9.1
%   Cn symmetry                        69.4x                          <- paper1_hpc 9.1
%   noise floor                        9.5-21.5% / 1.4% / 1.6-4.7%    <- paper1_hpc 6.1


\begin{document}

\twocolumn[
\begin{@twocolumnfalse}
\vspace*{-6mm}
\noindent{\small\bfseries\color{jpteal} Structural bioinformatics \quad Application Note\par}

\vspace{1mm}
\noindent{\LARGE\bfseries JuPIPER: editable scoring for FFT-based protein--protein docking on CPUs and GPUs\par}

\vspace{2.5mm}
\noindent{\normalsize Ilya Afanasyev$^{1,*}$\par}

\vspace{1.5mm}
\noindent{\footnotesize $^{1}$Innopolis University, Innopolis, Russia.
$^{*}$To whom correspondence should be addressed.
E-mail: \href{mailto:i.afanasyev@innopolis.ru}{i.afanasyev@innopolis.ru}%
\par}

\vspace{3mm}
\hrule
\vspace{2.5mm}

\noindent{\small\textbf{Abstract}\par}

\vspace{1mm}
\noindent\textbf{Summary}\quad
JuPIPER implements the published PIPER rigid-body docking search in Julia,
allowing researchers to change how shape and atomic contacts are scored
during the search on a CPU or an NVIDIA GPU. In bound redocking on one
RTX~4090, orientation searches for \campOk{} completed enzyme--inhibitor
cases from Docking Benchmark~4.0 took \campMinutes{} minutes in total, and in \campRankOne{} cases the lowest-energy placement was within
10\,\AA{} ligand C$\alpha$ RMSD of the deposited placement.

\smallskip
\noindent\textbf{Availability and Implementation}\quad
JuPIPER is implemented in Julia and distributed under the MIT licence.
The project repository is
\url{https://github.com/afanasyspb/JuPIPER}. The tagged software release
and its benchmark records are intended for archiving on Zenodo.
% Zenodo version DOI (after archiving):

\smallskip
\noindent\textbf{Contact}\quad
\href{mailto:i.afanasyev@innopolis.ru}{i.afanasyev@innopolis.ru}

\smallskip
\noindent\textbf{Supplementary Information}\quad
The supplement gives the correlation formula, CPU and GPU measurements,
a Hopf-fibre pruning test and results for all \campCases{} benchmark cases.
\vspace{2.5mm}\hrule\vspace{4mm}
\end{@twocolumnfalse}
]

\section{Introduction}
Protein--protein docking generates candidate structures of a complex
from structures of its two partners. The scoring function determines
which placements survive the initial search: a model discarded at this
stage is unavailable for later rescoring. Fast Fourier transforms (FFTs)
can score geometric fit at all translations for a given ligand orientation
(Katchalski-Katzir \textit{et al.}, 1992). PIPER adds interaction terms
(Kozakov \textit{et al.}, 2006) and underlies the ClusPro docking server
(Kozakov \textit{et al.}, 2017). JuPIPER implements this search in Julia
(Bezanson \textit{et al.}, 2017) for CPUs and NVIDIA GPUs. Method developers
can modify the physics-based FFT search or integrate it into a hybrid
pipeline without rewriting C or CUDA code.

For a structural biologist choosing interfaces to test, each pose predicts
which surfaces meet. Even at a fixed ligand orientation, changing the contact
weight can favour a different receptor face. JuPIPER exposes the weights,
intermediate fields and pose transformations behind that choice.

\section{Implementation}
The user supplies two structures and selects the receptor and ligand chains.
BioStructures reads the PDB coordinates. JuPIPER assigns atoms to 18 types, builds
the receptor fields and computes their Fourier transforms once. For each ligand
orientation on a Hopf grid (Yershova \textit{et al.}, 2010), the program rotates
the ligand, deposits its atoms on grids and evaluates all translations together
(Fig.~\ref{fig:campaign}a). The benchmark samples \campNrot{} orientations
at a translational grid spacing of 1.2\,\AA. Each orientation contributes
one pose, the translation with the lowest energy. Its rotation and translation
recover the ligand coordinates in the receptor frame.

\begin{figure*}[t]
\centering
\includegraphics[width=.96\textwidth]{fig1_pipeline.pdf}\par
\textbf{(a)}\par\smallskip
\includegraphics[width=.50\textwidth]{fig5_campaign_compact.pdf}\par
\textbf{(b)}\par
\caption{\textbf{(a)} The search reuses receptor fields across orientations
and retains the lowest-energy translation for each. Poses are
then ranked by energy and clustered by structural similarity.
\textbf{(b)} Bound-redocking recovery counts cases with at least one hit among
the $K$ lowest-energy poses. A hit has ligand C$\alpha$ RMSD $\leq 10$\,\AA{}
from the deposited placement with the receptor fixed. Recovery increases with
rank cut-off; the curves from the two hosts coincide.}
\label{fig:campaign}
\end{figure*}

The score combines shape, electrostatics and a contact potential. The benchmark
uses five components of the $18\times18$ DARS contact potential
(Chuang \textit{et al.}, 2008), giving seven fields in total. The shape
weight ($w_1=7$) sets the energy cost of placing ligand atoms in the receptor
core. The benchmark assigns zero atomic charges, so the electrostatic term
contributes zero even with $w_2=1500$. The contact weight ($w_3=2.5$)
scales the contribution of atom-type-specific contacts.

Clustering groups the 1,000 lowest-energy poses by pairwise ligand heavy-atom
root-mean-square deviation (RMSD), whereas benchmark recovery requires ligand
C$\alpha$ RMSD $\leq 10$\,\AA{} from the deposited pose, with the receptor fixed. Greedy
clustering with a 10\,\AA{} radius yields up to 30 centres, ranked by
the number of poses each represents. An energy rank therefore identifies
an individual placement, while a population rank identifies a family of
similar placements.

\subsection{Changing the contact weight}
From the repository root, start \texttt{julia -{}-project=test}.
This CPU example docks one chain from each PDB file on a small Hopf grid.
Replace the filenames and chain identifiers with those of your input structures.
Setting $w_3=3.0$ increases the benchmark contact weight by 20\% when
selecting the best translation for each orientation.
\begin{lstlisting}[basicstyle=\footnotesize\ttfamily,columns=fullflexible,
  keepspaces=true,showstringspaces=false,aboveskip=7pt,belowskip=7pt,
  frame=single,framerule=.3pt,framesep=3pt,xleftmargin=4pt,xrightmargin=4pt,
  rulecolor=\color{black!35},backgroundcolor=\color{black!3}]
using BioStructures, JuPIPER
include("examples/docking.jl")
rec = read_atoms("receptor.pdb"; chains="A")
lig = read_atoms("ligand.pdb"; chains="A")
rots = rotations_hopf_balanced(72)
poses, grid = search_poses(rec, lig, rots; w3=3.0)
best = argmin(p -> p.energy, poses)
R, t = pose_transform(best, rots, grid)
write_complex("best.pdb", rec, lig, R, t)
\end{lstlisting}
The output contains the receptor (chain R) and the lowest-energy ligand
placement (chain L) in the input receptor's coordinate frame. To examine
the effect of the weight alone, repeat the search on the same grid with
\code{w3=2.5} and superpose the two models on chain R.

On CPUs the correlation function uses FFTW (Frigo and Johnson, 2005).
On NVIDIA GPUs the same function calls cuFFT via CUDA.jl (Besard \textit{et al.}, 2019).
As in the Cartesian search described by Padhorny \textit{et al.}\ (2016),
weighted products are summed in Fourier space before a single inverse FFT.
The search reuses workspaces across ligand orientations.
Specialised routines handle accumulation and selection of the minimum.
Supplementary Figs.~S1(b) and S2 show device dispatch and the time spent
at each stage.

\section{Results}
\subsection{Effect of the shape weight}
For the bound E2A--HPr complex (1GGR), the published shape weight of 2.1 used
with standard van der Waals radii favoured placements inside the receptor.
A scan located the transition to repulsion at 3.151, close to the geometric
prediction of 3.264. The original PIPER description uses modified radii. With
standard radii and $w_1=7$, the lowest-energy placement had ligand C$\alpha$ RMSD
\dockRmsd\,\AA{} to the deposited ligand, measured with the receptor fixed.
The closest sampled pose reached \dockBest\,\AA{} at energy rank \dockBestRank.
On one RTX~4090, the \dockNrot{}-orientation search took \dockGpuFThirtyTwo\,s
in single precision. The full set of poses allows the effect of a rank cut-off
to be examined separately from the geometric resolution of the search.

\subsection{Bound redocking across the benchmark}
We used the enzyme--inhibitor cases of Docking Benchmark 4.0
(Hwang \textit{et al.}, 2010). Both partners were taken from their deposited
bound complex. The \campCases{} benchmark
cases represent 51 distinct structures because two use different chains of
1OYV. The control complex 1GGR was analysed separately.

Of the \campCases{} cases, \campOk{} completed. Atom typing failed for nine
atoms of a GLX residue in 4CPA. Among the completed cases, the lowest-energy
pose was a hit in \campRankOne{}/\campOk{} and the 1,000 lowest-energy poses
contained a hit in \campHitTopK{}/\campOk{} (Fig.~\ref{fig:campaign}b).
The centre of the most populated cluster was a hit in \campClusterOne{}/\campOk{}.
At the full depth of \campNrot{} poses, \campHitAll{}/\campOk{} had a hit.
The search for 1D6R completed without a hit and contributes to every denominator.
Across cases with a hit, the first-hit rank ranged from 1 to \campMaxHitRank{}.

Orientation searches for these cases took \campMinutes{} minutes in total on
one RTX~4090 in single precision (\campPerCaseLo--\campPerCaseHi\,s per case).
The timer excludes PDB input, field preparation, pose evaluation and clustering.
A repeat on another host with an RTX~4090 took \campTwoMinutes{} minutes.
First-hit rank, top-1,000 hit count and near-native cluster rank agreed for
\campEqualRanks{}/\campOk{} cases. The recorded lowest energies were identical
for \campEqualEnergy{}/\campOk{}.

\subsection{Comparison with ClusPro}
We submitted the first five targets in a predetermined case order to ClusPro.
The series size was fixed after the first comparison had been inspected.
The primary comparison required agreement on the population rank of the first
near-native cluster and a distance within the estimated joint discretisation
floor (1.6\,\AA{}). We measured the latter from the leading balanced ClusPro
model to the nearest JuPIPER cluster centre. The joint criterion was met in
\wpJoint{}/\wpCases{} cases. Both programs placed a near-native cluster first
in \wpSameBand{}/\wpCases{}. The nearest centre was as far as
\wpCentreAllHi\,\AA{} from the leading server model.

Among the 1,000 lowest-energy poses, \wpNearN{} of \wpCases{} searches
retained a pose within 2\,\AA{} of ClusPro's top-ranked balanced model
(Table~\ref{tab:wp8}). Those poses were \wpNearLo--\wpNearHi\,\AA{} away
at energy ranks \wpRankLo--\wpRankHi{}. For \wpOddCase{} the closest pose
was \wpOddNear\,\AA{} away. A user inspecting only the leading cluster
would miss this agreement at the level of individual poses.

\begin{table}[!ht]
\centering
\caption{Near-native centres are ranked by population (dash: none).
Ligand C$\alpha$ RMSD (\AA) to ClusPro's leading balanced model after
receptor alignment identifies the nearest JuPIPER centre (population rank)
and pose (energy rank). Individual poses more closely match the server model.}\label{tab:wp8}
\vspace{6pt}
\small\setlength{\tabcolsep}{2.7pt}
\begin{tabular}{@{}lrrrrrr@{}}\toprule
&\multicolumn{2}{c}{Near-native rank}&\multicolumn{2}{c}{Nearest centre}&\multicolumn{2}{c}{Nearest pose}\\
\cmidrule(lr){2-3}\cmidrule(lr){4-5}\cmidrule(l){6-7}
Case & ClusPro & JuPIPER & RMSD & Rank & RMSD & Rank\\\midrule
\wpTableBody
\bottomrule\end{tabular}
\end{table}

\section{Discussion}
Experimental evidence about contacting residues can guide the search in two
ways. Distance restraints exclude translations that separate specified sites
too far; an orientation can be skipped if none remain. Testing feasibility
still costs time. A residue mask instead modifies the receptor field to penalise
contacts outside a suspected interface, without changing the grid or rotation count.

In the unbound E2A--HPr test (1F3G and 1POH), requiring all 20 restraint groups to be satisfied
ruled out every translation for 78.06\% of orientations. Across the tested
cutoffs, restrained searches still took 1.26--1.86 times as long as the
unrestrained search. The near-native cluster grew from 34 to 435 poses and
rose from fifth to second by population, with interface RMSD remaining
within 2.37--2.50\,\AA{}. In a separate unbound test (1FGN and 1TFH),
doubling repulsion outside the antibody binding loops moved the centre
closest by interface RMSD from rank 17 to 15. Only in the E2A--HPr example
did the observed interface move into the two largest clusters. Neither
intervention sped up the search.

Bound redocking tests recovery from bound conformations, not predictive
performance across unbound targets. The five ClusPro comparisons do not
separate the effects of scoring, clustering and refinement. A matched runtime
comparison with PIPER remains open.

\section*{Data availability}
The benchmark case list, machine-tagged results, comparison JSON files and
figure generators accompany the source. The comparison tables can be regenerated
from the JSON files. Recomputing distances from coordinates also requires the
ClusPro model files. The repository records their checksums and job identifiers.
The service's two-month retention window for these submissions ended in November 2026.

\section*{Conflict of interest}
None declared.

\begin{thebibliography}{}
\small\setlength{\itemsep}{1.5pt}
\item Besard T, Foket C, De Sutter B. Effective extensible programming:
unleashing Julia on GPUs. \textit{IEEE Trans Parallel Distrib Syst}
2019; \textbf{30}: 827--41.

\item Bezanson J, Edelman A, Karpinski S \textit{et al.} Julia: a fresh approach
to numerical computing. \textit{SIAM Rev} 2017; \textbf{59}: 65--98.

\item Chuang G-Y, Kozakov D, Brenke R \textit{et al.} DARS (Decoys As the
Reference State) potentials for protein--protein docking. \textit{Biophys J}
2008; \textbf{95}: 4217--27.

\item Frigo M, Johnson SG. The design and implementation of FFTW3.
\textit{Proc IEEE} 2005; \textbf{93}: 216--31.

\item Hwang H, Vreven T, Janin J \textit{et al.} Protein--protein docking
benchmark version 4.0. \textit{Proteins} 2010; \textbf{78}: 3111--4.

\item Katchalski-Katzir E, Shariv I, Eisenstein M \textit{et al.} Molecular
surface recognition: determination of geometric fit between proteins and their
ligands by correlation techniques. \textit{Proc Natl Acad Sci USA}
1992; \textbf{89}: 2195--9.

\item Kozakov D, Brenke R, Comeau SR \textit{et al.} PIPER: an FFT-based protein
docking program with pairwise potentials. \textit{Proteins} 2006; \textbf{65}:
392--406.

\item Kozakov D, Hall DR, Xia B \textit{et al.} The ClusPro web server for
protein--protein docking. \textit{Nat Protoc} 2017; \textbf{12}: 255--78.

\item Padhorny D, Kazennov A, Zerbe BS \textit{et al.} Protein--protein docking
by fast generalized Fourier transforms on 5D rotational manifolds.
\textit{Proc Natl Acad Sci USA} 2016; \textbf{113}: E4286--93.

\item Yershova A, Jain S, LaValle SM \textit{et al.} Generating uniform
incremental grids on SO(3) using the Hopf fibration. \textit{Int J Robot Res}
2010; \textbf{29}: 801--12.
\end{thebibliography}
\end{document}
